\documentclass[11pt]{article}
\usepackage[margin=0.8in]{geometry}
\usepackage{amsmath,amssymb,amsthm,hyperref}
\hypersetup{colorlinks=true,urlcolor=blue,linkcolor=blue}
\newtheorem{theorem}{Theorem}
\newtheorem{lemma}{Lemma}
\newcommand{\CF}[2]{\mathop{\mathrm K}_{j=1}^{#1}\frac{#2_j}{x}}
\title{Identifying the telescoping continued-fraction coefficients\\
\large MathDB 354157: the stabilized coefficients equal the Stirling coefficients}
\author{Principia Math In Progress}
\date{12 September 2026}
\begin{document}
\maketitle

\section*{Task 1. Prove coefficient matching}
Bhatnagar and Rajkumar's Conjecture 6.2 in \cite{BR}, p.~9,
asks whether the coefficients obtained by their telescoping algorithm
agree with those of the classical continued fraction for the Stirling
remainder. The algorithm first fixes a Legendre approximation order
\(k\), so denote its output by \(a_j^{(k)}\).
The sequence in the conjecture is the stabilized sequence described
on p.~7 and displayed in (6.1). We prove its stabilization and
identification together:

\begin{theorem}\label{main}
Let \(b_j\) be the classical Stirling continued-fraction coefficients.
For every \(m\geq1\) and \(k\geq m+1\), the source's successive
leading-degree cancellations are possible through step \(m\) and give
\[
 a_j^{(k)}=b_j\qquad(1\leq j\leq m).
\]
Consequently the stabilized sequence \(a_j\) satisfies \(a_j=b_j\)
for every \(j\), resolving Conjecture 6.2.
\end{theorem}

The finite-\(k\) convention matters: for example,
\(a_2^{(2)}=1/18\), whereas the stabilized value is \(b_2=1/30\).
The theorem concerns exactly the stabilization region \(k\geq j+1\)
specified in the source.

\subsection*{The two formal expansions}
Let
\[
 r(x)=\log\Gamma(x+1)-(x+\tfrac12)\log x+x-\tfrac12\log(2\pi),
 \qquad x>0.
\]
Its formal Stirling expansion and exact difference are
\begin{align}
 \mathcal R(x)&=\sum_{\ell\geq1}
 \frac{B_{2\ell}}{2\ell(2\ell-1)x^{2\ell-1}},
                                                        \label{Stirling}\\
 r(x)-r(x+1)&=\varepsilon(x):=
 (x+\tfrac12)\log(1+1/x)-1.                              \label{difference}
\end{align}
Thus the formal expansion of \(\varepsilon\) is
\(\mathcal R(x)-\mathcal R(x+1)\).
All formal \(O\)-statements below refer to Laurent expansions at infinity.

The classical coefficients \(b_j>0\) appearing in the conjecture
are characterized by the formal continued fraction
\[
 \mathcal R(x)=\cfrac{b_1}{x+\cfrac{b_2}{x+\cfrac{b_3}{x+\ddots}}}.
\]
In particular, if \(B_m(x)=\CF{m}{b}\), then
\begin{equation}\label{formal}
 \mathcal R(x)-B_m(x)=O(x^{-2m-1}).
\end{equation}
These are the defining formal coefficients of the classical
Stieltjes expansion recalled in \cite[Section 6]{BR}.
Only this formal property and nonvanishing of \(b_j\) are needed.

For completeness, the positive Stieltjes measure underlying these
facts follows by integrating Binet's formula by parts:
\[
 \frac{r(\sqrt z)}{\sqrt z}
 =\int_0^\infty\frac{\rho(t)}{z+t}\,dt,\qquad
 \rho(t)=\frac{-\log(1-e^{-2\pi\sqrt t})}{2\pi\sqrt t}>0.
\]
Its moments are
\(\mu_n=|B_{2n+2}|/[(2n+2)(2n+1)]\).
Every ordinary and shifted Hankel moment determinant is positive,
because its matrix is the Gram matrix of monomials for respectively
\(\rho(t)\,dt\) and \(t\rho(t)\,dt\), both with infinite support.
The classical Stieltjes correspondence therefore gives strictly
positive, nonterminating continued-fraction coefficients.
No convergence assertion is used to infer \eqref{formal}.

\subsection*{Accuracy of the finite Legendre input}
Let \(P_k\) be the Legendre polynomial and \(P_k^*\) its associated
numerator polynomial, with the normalization of \cite[Section 2]{BR}.
Put
\[
 E_k(x)=\frac{(2x+1)P_k^*(2x+1)}{P_k(2x+1)}-1
       =\frac{f_k^*(z)}{f_k(z)},\qquad z=x(x+1).
\]
For \(k\geq2\), \(\deg f_k=d_k\) and \(\deg f_k^*=d_k-1\), after the
usual cancellation when \(k\) is odd.

\begin{lemma}\label{Legendre}
For each \(k\geq1\),
\[
 \varepsilon(x)-E_k(x)=O(x^{-2k}).
\]
\end{lemma}
\begin{proof}
The associated numerator has the integral expression
\[
 P_k^*(t)=\frac12\int_{-1}^1\frac{P_k(t)-P_k(s)}{t-s}\,ds.
\]
Consequently
\[
 \frac12\log\frac{t+1}{t-1}-\frac{P_k^*(t)}{P_k(t)}
 =\frac1{2P_k(t)}\int_{-1}^1\frac{P_k(s)}{t-s}\,ds.
\]
Expanding \(1/(t-s)\) at infinity, orthogonality of \(P_k\)
to \(1,s,\ldots,s^{k-1}\) makes the integral \(O(t^{-k-1})\).
Division by the degree-\(k\) polynomial makes the right side
\(O(t^{-2k-1})\). Multiplying by \(t\) and setting \(t=2x+1\)
gives the assertion.
\end{proof}

\subsection*{Leading-degree cancellation identifies the next coefficient}
For any prefix \(a_1,\ldots,a_m\), write
\[
 G_m(x)=\CF{m}{a}=\frac{p_m(x)}{q_m(x)}.
\]
The continuant recurrence is
\[
 p_m=xp_{m-1}+a_mp_{m-2},\qquad
 q_m=xq_{m-1}+a_mq_{m-2},
\]
with \((p_{-1},p_0)=(1,0)\) and \((q_{-1},q_0)=(0,1)\).
In particular \(q_m\) is monic of degree \(m\), and
\begin{equation}\label{continuant}
 G_m-G_{m-1}
 =\frac{(-1)^{m-1}a_1\cdots a_m}{q_mq_{m-1}}.
\end{equation}
Parity of the continuants shows that \(G_m\) is odd.
Hence \(G_m(x)-G_m(x+1)\) is invariant under \(x\mapsto-x-1\);
its numerator and denominator in the source's representation
are polynomials in \(z=x(x+1)\).

Fix \(k\geq m+1\), and suppose the first \(m-1\) chosen coefficients
are \(b_1,\ldots,b_{m-1}\). Leave the final coefficient \(a\) free
and write \(G_m(x;a)\). Formula \eqref{continuant} gives
\begin{equation}\label{variation}
 G_m(x;a)-B_m(x)
 =(-1)^{m-1}\!\left(\prod_{j=1}^{m-1}b_j\right)
 (a-b_m)x^{-2m+1}+O(x^{-2m-1}).
\end{equation}
For any \(q\geq1\), a leading term \(c x^{-q}\) has difference
\[
 c x^{-q}-c(x+1)^{-q}=q c x^{-q-1}+O(x^{-q-2}).
\]
Moreover, Lemma \ref{Legendre}, \eqref{difference} and
\eqref{formal} show that
\begin{equation}\label{good}
 E_k(x)-\bigl(B_m(x)-B_m(x+1)\bigr)=O(x^{-2m-2}),
\end{equation}
since \(k\geq m+1\). Combining \eqref{variation} and \eqref{good}
therefore yields
\begin{align}
 E_k(x)-\bigl(G_m(x;a)-G_m(x+1;a)\bigr)
 ={}&(2m-1)(-1)^m
 \left(\prod_{j=1}^{m-1}b_j\right)(a-b_m)x^{-2m}\nonumber\\
 &+O(x^{-2m-1}).                                        \label{leading}
\end{align}
The coefficient of \(a-b_m\) is nonzero.

Now relate this directly to the polynomial algorithm.
With \(q_m^+(x)=q_m(x+1)\) and similarly for \(p_m^+\), the
cleared numerator is exactly
\[
 \Delta_1(k,m)=f_k^*q_mq_m^+
              -f_k(p_mq_m^+-p_m^+q_m).
\]
Its denominator is \(f_k(z)q_m(x)q_m(x+1)\).
As a polynomial in \(z\), this denominator has degree \(d_k+m\)
and a fixed nonzero leading coefficient. The order estimate
\eqref{leading} thus says that \(\Delta_1(k,m)\) has degree at
most \(d_k\), and its coefficient of \(z^{d_k}\) is a fixed
nonzero multiple of \(a-b_m\).
The source's rule is precisely to set this highest possible
coefficient equal to zero. Its unique choice is therefore \(a=b_m\).
For this choice, the stronger estimate \eqref{good} gives
\(\deg\Delta_1(k,m)\leq d_k-1\), as required.

Starting with the empty prefix proves the assertion inductively.
There is no premature termination before the required last step:
after a step \(j<m\), \eqref{continuant} and \eqref{formal} give
\[
 \mathcal R(x)-B_j(x)
 =(-1)^j\!\left(\prod_{i=1}^{j+1}b_i\right)x^{-2j-1}
  +O(x^{-2j-3}).
\]
Because \(k\geq m+1\geq j+2\), the Legendre-input error is of
strictly higher order than the difference of this expression.
Thus the residual has a nonzero term of order \(x^{-2j-2}\);
its cleared numerator has degree exactly \(d_k-1\).
Every next cancellation above is well defined. This proves
Theorem \ref{main}.\qed

\subsection*{Consequences and exact verification}
The full stabilized continued fraction is therefore the classical
Stirling one. Its standard Stieltjes bounds, already recalled in
\cite[Section 6]{BR}, have even convergents below \(r(x)\) and odd
convergents above \(r(x)\), for \(x>0\).
This also gives the intended companion bound in MathDB 354156;
no separate campaign credit is taken for that immediate consequence.
The source's displayed parity inequalities in Conjecture 6.1 are
reversed relative to its prose and preceding results; we use the
unambiguous intended parity.

The verifier \texttt{verify\_coefficients.py} implements two independent
exact constructions. One builds Legendre polynomials, continuants
and cleared residuals, then cancels their leading coefficients.
The other starts with the Bernoulli series and repeatedly extracts
the continued fraction using formal-series reciprocation.
All \(45\) coefficients in the prefixes \(2\leq k\leq10\),
\(1\leq m<k\), agree. It separately checks the nonstabilized
value \(a_2^{(2)}=1/18\). The all-order proof above, not these
finite checks, establishes the conjecture.

\section*{Tasks 2--5. Further directions after the full solve}
These are proposed tasks. Check their current literature and public
claims before attacking them.
\begin{description}
\item[Task 2.] For fixed \(k\), determine whether all coefficients
\(a_m^{(k)}\) beyond the stabilization region exist and remain positive.
Identify the asymptotic behavior of this finite-\(k\) tail.
\item[Task 3.] Settle the source's separate sign conjecture for all
coefficients of the cleared residual \(\Delta_1(k,m)\) after its
leading-degree cancellation.
\item[Task 4.] Derive explicit convergence-rate bounds for the
stabilized convergents \(B_m(x)\), uniform on \(x\geq x_0>0\),
with usable dependence on both \(m\) and \(x_0\).
\item[Task 5.] Characterize positive moment sequences whose formal
discrete antiderivatives also admit positive Stieltjes
continued fractions, extending the structural mechanism behind
the Stirling example.
\end{description}

\begin{thebibliography}{9}
\bibitem{BR}
G. Bhatnagar and K. Rajkumar,
\emph{Telescoping continued fractions for the error term in Stirling's formula},
J. Approx. Theory 293 (2023), 105943,
\url{https://doi.org/10.1016/j.jat.2023.105943}.
Current author-linked manuscript:
\url{https://arxiv.org/abs/2204.00962v2}.
Earlier version, with the stabilization results in Sections 6--7:
\url{https://arxiv.org/abs/2204.00962v1}.
\end{thebibliography}
\end{document}
